\section{Discussion}

\subsection{Implications for Review Process Design}

Our findings suggest specific design recommendations:
\begin{itemize}
    \item \textbf{Two rounds are usually sufficient}: 85\% of papers reach threshold within 2 rounds. Budget for 2 rounds; only extend to 3+ for papers with many P1 items.
    \item \textbf{P1 triage is critical}: The synthesis step's P1/P2/P3 classification focuses author effort where it matters most. Without triage, authors may spend time on P3 items while P1 issues remain.
    \item \textbf{Initial score predicts effort}: Papers starting below 5.0/10 require significantly more revision than those starting above 6.0/10.
\end{itemize}

\subsection{Limitations}

\textbf{AI-simulated vs.\ real reviews}. Score trajectories may differ with real venue reviewers, who bring deeper domain knowledge and different biases. The absolute scores should be interpreted as relative quality indicators, not predictions of venue outcomes.

\textbf{Same-author confound}. All 14 papers are by the same author, meaning revision quality is confounded with a single author's revision skill. Multi-author validation is needed.

\textbf{Score inflation risk}. Reviewers in later rounds may be primed by seeing improvements, leading to inflated scores. We mitigate this by generating round N reviews independently, but the risk remains.

\subsection{Comparison to Real Venue Dynamics}

Published data on NeurIPS rebuttal outcomes shows that approximately 20--30\% of papers change their accept/reject status after rebuttal. Our 85\% improvement rate is higher, likely because: (1) authors implement full revisions rather than just rebuttals, and (2) AI reviewers may be more responsive to addressed concerns.
